\section{Phản tư}

Quá trình nghiên cứu và thực nghiệm cải tiến hệ thống truyền thông tại daotao.ai đã mang lại cho tôi những bài học thực tế dưới vai trò là một chuyên viên vận hành trực tiếp tại Trung tâm Công nghệ Giáo dục.

\subsection{Sự thay đổi nhận thức về khái niệm truyền thông}

Trước khi thực hiện bài luận này, tôi từng có góc nhìn đơn giản về truyền thông trong đào tạo trực tuyến, xem đó chỉ là hoạt động gửi thông tin kỹ thuật một chiều như soạn thông báo học tập và gửi đi. Tuy nhiên, dưới lăng kính của mô hình giao dịch \parencite{barnlund2008} và khái niệm trường kinh nghiệm \parencite{schramm1954}, tôi nhận ra truyền thông thực chất là quá trình kiến tạo sự hiểu biết chung giữa các bên.

Việc người học không kích hoạt được tài khoản hay hiểu sai thông số không hoàn toàn do lỗi của họ, mà do người vận hành chưa thấu hiểu trường kinh nghiệm công nghệ của người nhận là các học viên lớn tuổi thuộc Đề án 06. Truyền thông thành công đòi hỏi người thiết kế hệ thống phải chuyển dịch thông điệp từ thuật ngữ kỹ thuật sang ngôn từ trực quan, gần gũi với người học.

\subsection{Vai trò của công nghệ trong đào tạo trực tuyến}

Nghiên cứu giúp tôi nhận thức rõ hơn về vai trò của công nghệ trong truyền thông giáo dục trực tuyến:
\begin{itemize}
    \item \textbf{Về mặt tích cực:} Công nghệ là công cụ hỗ trợ vận hành hiệu quả thông qua tự động hóa, thể hiện qua việc chatbot giúp xử lý phần lớn yêu cầu lặp lại.
    \item \textbf{Về mặt hạn chế:} Nếu công nghệ không được thiết kế lấy người học làm trung tâm, nó sẽ tạo ra rào cản kỹ thuật và đặt ra gánh nặng tải nhận thức ngoại lai cho học viên \parencite{sweller1988}.
\end{itemize}

Một nút đăng ký đặt sai vị trí hay một email hướng dẫn nhiều chữ kỹ thuật chính là minh chứng cho thấy công nghệ có thể tạo ra rào cản nhận thức nếu không được thiết kế hợp lý.

\subsection{Giới hạn giải pháp và bài học thực tiễn}

Tôi ý thức được các giới hạn của nghiên cứu này: chatbot hiện tại mới chỉ giải quyết các lỗi phổ biến ở dạng đóng, chưa có khả năng giải quyết các lỗi cá biệt hoặc lỗi phát sinh ngoài kịch bản. Mẫu email hướng dẫn cũng cần phải được kiểm thử liên tục với các tập mẫu lớn hơn để tối ưu cấu trúc nội dung.

Bài học lớn nhất mà tôi rút ra cho công việc vận hành là: dạy học trực tuyến thực chất là thiết kế truyền thông có mục tiêu. Một chuyên viên vận hành cần có kỹ năng thiết kế trải nghiệm giao tiếp số, thấu hiểu khó khăn của người học và kiên trì giảm thiểu các rào cản truyền thông nhằm giúp học viên có được hành trình học tập thuận lợi.
